% Options for packages loaded elsewhere
\PassOptionsToPackage{unicode}{hyperref}
\PassOptionsToPackage{hyphens}{url}
\PassOptionsToPackage{dvipsnames,svgnames,x11names}{xcolor}
%
\documentclass[
  11pt,
]{article}
\usepackage{amsmath,amssymb}
\usepackage{iftex}
\ifPDFTeX
  \usepackage[T1]{fontenc}
  \usepackage[utf8]{inputenc}
  \usepackage{textcomp} % provide euro and other symbols
\else % if luatex or xetex
  \usepackage{unicode-math} % this also loads fontspec
  \defaultfontfeatures{Scale=MatchLowercase}
  \defaultfontfeatures[\rmfamily]{Ligatures=TeX,Scale=1}
\fi
\usepackage{lmodern}
\ifPDFTeX\else
  % xetex/luatex font selection
\fi
% Use upquote if available, for straight quotes in verbatim environments
\IfFileExists{upquote.sty}{\usepackage{upquote}}{}
\IfFileExists{microtype.sty}{% use microtype if available
  \usepackage[]{microtype}
  \UseMicrotypeSet[protrusion]{basicmath} % disable protrusion for tt fonts
}{}
\makeatletter
\@ifundefined{KOMAClassName}{% if non-KOMA class
  \IfFileExists{parskip.sty}{%
    \usepackage{parskip}
  }{% else
    \setlength{\parindent}{0pt}
    \setlength{\parskip}{6pt plus 2pt minus 1pt}}
}{% if KOMA class
  \KOMAoptions{parskip=half}}
\makeatother
\usepackage{xcolor}
\usepackage[margin=1in]{geometry}
\usepackage{color}
\usepackage{fancyvrb}
\newcommand{\VerbBar}{|}
\newcommand{\VERB}{\Verb[commandchars=\\\{\}]}
\DefineVerbatimEnvironment{Highlighting}{Verbatim}{commandchars=\\\{\}}
% Add ',fontsize=\small' for more characters per line
\usepackage{framed}
\definecolor{shadecolor}{RGB}{248,248,248}
\newenvironment{Shaded}{\begin{snugshade}}{\end{snugshade}}
\newcommand{\AlertTok}[1]{\textcolor[rgb]{0.94,0.16,0.16}{#1}}
\newcommand{\AnnotationTok}[1]{\textcolor[rgb]{0.56,0.35,0.01}{\textbf{\textit{#1}}}}
\newcommand{\AttributeTok}[1]{\textcolor[rgb]{0.13,0.29,0.53}{#1}}
\newcommand{\BaseNTok}[1]{\textcolor[rgb]{0.00,0.00,0.81}{#1}}
\newcommand{\BuiltInTok}[1]{#1}
\newcommand{\CharTok}[1]{\textcolor[rgb]{0.31,0.60,0.02}{#1}}
\newcommand{\CommentTok}[1]{\textcolor[rgb]{0.56,0.35,0.01}{\textit{#1}}}
\newcommand{\CommentVarTok}[1]{\textcolor[rgb]{0.56,0.35,0.01}{\textbf{\textit{#1}}}}
\newcommand{\ConstantTok}[1]{\textcolor[rgb]{0.56,0.35,0.01}{#1}}
\newcommand{\ControlFlowTok}[1]{\textcolor[rgb]{0.13,0.29,0.53}{\textbf{#1}}}
\newcommand{\DataTypeTok}[1]{\textcolor[rgb]{0.13,0.29,0.53}{#1}}
\newcommand{\DecValTok}[1]{\textcolor[rgb]{0.00,0.00,0.81}{#1}}
\newcommand{\DocumentationTok}[1]{\textcolor[rgb]{0.56,0.35,0.01}{\textbf{\textit{#1}}}}
\newcommand{\ErrorTok}[1]{\textcolor[rgb]{0.64,0.00,0.00}{\textbf{#1}}}
\newcommand{\ExtensionTok}[1]{#1}
\newcommand{\FloatTok}[1]{\textcolor[rgb]{0.00,0.00,0.81}{#1}}
\newcommand{\FunctionTok}[1]{\textcolor[rgb]{0.13,0.29,0.53}{\textbf{#1}}}
\newcommand{\ImportTok}[1]{#1}
\newcommand{\InformationTok}[1]{\textcolor[rgb]{0.56,0.35,0.01}{\textbf{\textit{#1}}}}
\newcommand{\KeywordTok}[1]{\textcolor[rgb]{0.13,0.29,0.53}{\textbf{#1}}}
\newcommand{\NormalTok}[1]{#1}
\newcommand{\OperatorTok}[1]{\textcolor[rgb]{0.81,0.36,0.00}{\textbf{#1}}}
\newcommand{\OtherTok}[1]{\textcolor[rgb]{0.56,0.35,0.01}{#1}}
\newcommand{\PreprocessorTok}[1]{\textcolor[rgb]{0.56,0.35,0.01}{\textit{#1}}}
\newcommand{\RegionMarkerTok}[1]{#1}
\newcommand{\SpecialCharTok}[1]{\textcolor[rgb]{0.81,0.36,0.00}{\textbf{#1}}}
\newcommand{\SpecialStringTok}[1]{\textcolor[rgb]{0.31,0.60,0.02}{#1}}
\newcommand{\StringTok}[1]{\textcolor[rgb]{0.31,0.60,0.02}{#1}}
\newcommand{\VariableTok}[1]{\textcolor[rgb]{0.00,0.00,0.00}{#1}}
\newcommand{\VerbatimStringTok}[1]{\textcolor[rgb]{0.31,0.60,0.02}{#1}}
\newcommand{\WarningTok}[1]{\textcolor[rgb]{0.56,0.35,0.01}{\textbf{\textit{#1}}}}
\usepackage{longtable,booktabs,array}
\usepackage{calc} % for calculating minipage widths
% Correct order of tables after \paragraph or \subparagraph
\usepackage{etoolbox}
\makeatletter
\patchcmd\longtable{\par}{\if@noskipsec\mbox{}\fi\par}{}{}
\makeatother
% Allow footnotes in longtable head/foot
\IfFileExists{footnotehyper.sty}{\usepackage{footnotehyper}}{\usepackage{footnote}}
\makesavenoteenv{longtable}
\setlength{\emergencystretch}{3em} % prevent overfull lines
\providecommand{\tightlist}{%
  \setlength{\itemsep}{0pt}\setlength{\parskip}{0pt}}
\setcounter{secnumdepth}{-\maxdimen} % remove section numbering
\ifLuaTeX
  \usepackage{selnolig}  % disable illegal ligatures
\fi
\IfFileExists{bookmark.sty}{\usepackage{bookmark}}{\usepackage{hyperref}}
\IfFileExists{xurl.sty}{\usepackage{xurl}}{} % add URL line breaks if available
\urlstyle{same}
\hypersetup{
  pdftitle={Posterior Uncertainty in the Whitened SVD Basis},
  pdfauthor={Christian Hellum Bye},
  colorlinks=true,
  linkcolor={NavyBlue},
  filecolor={Maroon},
  citecolor={Blue},
  urlcolor={NavyBlue},
  pdfcreator={LaTeX via pandoc}}

\title{Posterior Uncertainty in the Whitened SVD Basis}
\usepackage{etoolbox}
\makeatletter
\providecommand{\subtitle}[1]{% add subtitle to \maketitle
  \apptocmd{\@title}{\par {\large #1 \par}}{}{}
}
\makeatother
\subtitle{mistsim map-making: from the SVD of \(\tilde{A}\) to
\(\sigma^2_{\rm post}/\sigma^2_{\rm prior}\)}
\author{Christian Hellum Bye}
\date{11 September 2026}

\begin{document}
\maketitle

{
\hypersetup{linkcolor=}
\setcounter{tocdepth}{2}
\tableofcontents
}
How \texttt{mistsim} turns an SVD of the whitened design matrix into a
posterior standard-deviation map, and what the
\(\sigma^2_{\rm post}/\sigma^2_{\rm prior}\) plots actually show.

\textbf{Code:}

\begin{longtable}[]{@{}ll@{}}
\toprule\noalign{}
Piece & Location \\
\midrule\noalign{}
\endhead
\bottomrule\noalign{}
\endlastfoot
Whitened design matrix \(\tilde{A}\) &
\texttt{src/mistsim/mapmaking.py:281} (\texttt{make\_Atilde}) \\
Prior diagonal \(S\) & \texttt{src/mistsim/pipeline.py:565}
(\texttt{compute\_prior}) \\
Truncated SVD & \texttt{src/mistsim/pipeline.py:588}
(\texttt{run\_svd}) \\
Mode selection & \texttt{src/mistsim/pipeline.py:604}
(\texttt{select\_nvec}) \\
Wiener mean & \texttt{src/mistsim/pipeline.py:749}
(\texttt{wiener\_filter}) \\
Posterior sampling & \texttt{src/mistsim/pipeline.py:1026}
(\texttt{posterior\_uncertainty}) \\
Plotting & \texttt{src/mistsim/plotting.py:949}
(\texttt{plot\_posterior\_maps}) \\
\end{longtable}

Rendered examples live in \texttt{notebooks/mapmaking/notebooks\_out/}
(cell 21 of each executed notebook).

\begin{center}\rule{0.5\linewidth}{0.5pt}\end{center}

\hypertarget{1-the-model}{%
\subsection{1. The model}\label{1-the-model}}

The forward problem is

\[ y = A x + n \]

where \(x\) is the packed real spherical-harmonic vector, \(A\) the
beam-convolution / drift-scan operator (\texttt{make\_Amat},
\texttt{mapmaking.py:188}), \(y\) the simulated timestream, and \(n\)
the radiometer noise.

Two Gaussian ingredients close the problem:

\begin{itemize}
\tightlist
\item
  \textbf{Noise:} \(n \sim \mathcal{N}(0, N)\) with
  \(N = \mathrm{diag}(\texttt{Ndiag})\), one variance per (time,
  frequency) sample.
\item
  \textbf{Prior:} \(x \sim \mathcal{N}(0, S)\) with
  \(S = \mathrm{diag}(\texttt{Sdiag})\), where
  \texttt{Sdiag\ =\ cl{[}ells\_full{]}} --- the sky power spectrum
  evaluated at each coefficient\textquotesingle s \(\ell\).
\end{itemize}

\hypertarget{why-the-prior-is-exactly-diagonal}{%
\subsubsection{Why the prior is exactly
diagonal}\label{why-the-prior-is-exactly-diagonal}}

\texttt{compute\_prior} assigns \(C_\ell\) to \emph{every} real degree
of freedom, including the real and imaginary parts of \(m > 0\)
separately:

\begin{Shaded}
\begin{Highlighting}[]
\NormalTok{ells\_pos  }\OperatorTok{=}\NormalTok{ ells\_hp[emms\_hp }\OperatorTok{!=} \DecValTok{0}\NormalTok{]}
\NormalTok{ells\_full }\OperatorTok{=}\NormalTok{ np.concatenate((ells\_hp, ells\_pos))}
\ControlFlowTok{return}\NormalTok{ cl[ells\_full]}
\end{Highlighting}
\end{Shaded}

This is consistent because of the \(1/\sqrt{2}\) in the packing
convention (\texttt{alm.py:122}):

\[ a_{\ell m} = \tfrac{1}{\sqrt{2}}\left(a^{\rm re}_{\ell m} + i\,a^{\rm im}_{\ell m}\right) \]

A complex coefficient with \(\langle |a_{\ell m}|^2 \rangle = C_\ell\)
splits into two real d.o.f. each with variance \(C_\ell\). So in the
packed basis the prior covariance really is diagonal with entries
\(C_\ell\) --- no \(m\)-dependent bookkeeping, and the whitened prior
below is exactly the identity.

\begin{center}\rule{0.5\linewidth}{0.5pt}\end{center}

\hypertarget{2-whitening}{%
\subsection{2. Whitening}\label{2-whitening}}

Define

\[ \tilde{x} = S^{-1/2} x, \qquad \tilde{y} = N^{-1/2} y, \qquad \tilde{A} = N^{-1/2} A\, S^{1/2} \]

so that

\[ \tilde{y} = \tilde{A}\,\tilde{x} + \tilde{n}, \qquad \tilde{x} \sim \mathcal{N}(0, I), \quad \tilde{n} \sim \mathcal{N}(0, I). \]

This is \texttt{make\_Atilde} (\texttt{mapmaking.py:281}), applied
matrix-free:

\begin{Shaded}
\begin{Highlighting}[]
\KeywordTok{def}\NormalTok{ \_Atilde\_matvec(v, Ndiag, Amat, Sdiag):}
\NormalTok{    Nm12 }\OperatorTok{=} \DecValTok{1} \OperatorTok{/}\NormalTok{ np.sqrt(Ndiag)}
\NormalTok{    S12 }\OperatorTok{=}\NormalTok{ np.sqrt(Sdiag)}
    \ControlFlowTok{return}\NormalTok{ Nm12 }\OperatorTok{*}\NormalTok{ Amat.matvec(S12 }\OperatorTok{*}\NormalTok{ v)}
\end{Highlighting}
\end{Shaded}

and on the data side,
\texttt{y\_tilde\ =\ Ndiag**(-0.5)\ *\ (y\ +\ noise)}
(\texttt{pipeline.py:776}).

\textbf{Why bother.} After whitening, the prior and the noise are the
same object --- the identity. A single SVD of \(\tilde{A}\) therefore
diagonalizes the entire posterior, and the singular values come out as
pure dimensionless signal-to-noise numbers: \(\sigma_k\) is how well the
data constrains mode \(k\) \emph{in units of its own prior width}.
Nothing carries units, and no mode is privileged by an arbitrary choice
of amplitude scale.

\begin{center}\rule{0.5\linewidth}{0.5pt}\end{center}

\hypertarget{3-the-posterior}{%
\subsection{3. The posterior}\label{3-the-posterior}}

With both covariances the identity,

\[ -2 \ln p(\tilde{x} \mid \tilde{y}) = \lVert \tilde{y} - \tilde{A}\tilde{x} \rVert^2 + \lVert \tilde{x} \rVert^2 + \text{const} \]

Both terms are quadratic, so the posterior is Gaussian:

\[ C_{\rm post} = \left(\tilde{A}^{\mathsf H}\tilde{A} + I\right)^{-1}, \qquad \mu = C_{\rm post}\,\tilde{A}^{\mathsf H} \tilde{y} \]

The \(+I\) \emph{is} the prior. Drop it and this is ordinary least
squares, which is singular: the drift scan has exact null directions
(modes the beam never sees), and those blow up.

\begin{center}\rule{0.5\linewidth}{0.5pt}\end{center}

\hypertarget{4-the-svd-diagonalizes-it}{%
\subsection{4. The SVD diagonalizes
it}\label{4-the-svd-diagonalizes-it}}

Write \(\tilde{A} = U \Sigma V^{\mathsf H}\) (\texttt{run\_svd},
\texttt{pipeline.py:588}, via \texttt{scipy.sparse.linalg.svds}). Then

\[ \tilde{A}^{\mathsf H}\tilde{A} + I = V\left(\Sigma^2 + I\right)V^{\mathsf H} \quad\Longrightarrow\quad C_{\rm post} = V\left(\Sigma^2 + I\right)^{-1}V^{\mathsf H} \]

The right singular vectors \(v_k\) are the principal axes of the
posterior. Along \(v_k\):

\begin{longtable}[]{@{}ll@{}}
\toprule\noalign{}
& value \\
\midrule\noalign{}
\endhead
\bottomrule\noalign{}
\endlastfoot
prior std (whitened) & \(1\) \\
posterior std & \(1/\sqrt{\sigma_k^2 + 1}\) \\
variance ratio & \(1/(\sigma_k^2 + 1)\) \\
\end{longtable}

In code (\texttt{pipeline.py:1052}):

\begin{Shaded}
\begin{Highlighting}[]
\NormalTok{post\_std\_svd }\OperatorTok{=} \FloatTok{1.0} \OperatorTok{/}\NormalTok{ np.sqrt(Dnum}\OperatorTok{**}\DecValTok{2} \OperatorTok{+} \FloatTok{1.0}\NormalTok{)}
\NormalTok{std\_reduce }\OperatorTok{=} \DecValTok{1} \OperatorTok{{-}}\NormalTok{ post\_std\_svd}
\end{Highlighting}
\end{Shaded}

The two limits are the whole story:

\begin{itemize}
\tightlist
\item
  \(\sigma_k \gg 1\) --- well-measured mode. Posterior std
  \(\to 1/\sigma_k\), far tighter than the prior. The data dominates.
\item
  \(\sigma_k \ll 1\) --- unconstrained mode. Posterior std \(\to 1\):
  the posterior hands back the prior unchanged, which is the correct and
  honest answer for something the instrument never saw.
\end{itemize}

\texttt{std\_reduce} is the \emph{shrinkage} --- the fraction of the
prior width that the data removes.

\begin{center}\rule{0.5\linewidth}{0.5pt}\end{center}

\hypertarget{5-the-eigenmodes-of-the-posterior}{%
\subsection{5. The eigenmodes of the
posterior}\label{5-the-eigenmodes-of-the-posterior}}

The SVD has already produced them. Section 4 wrote

\[ \tilde{C}_{\rm post} = V\left(\Sigma^2+I\right)^{-1}V^{\mathsf H} + \left(I - VV^{\mathsf H}\right) \]

which \emph{is} an eigendecomposition: the rows of \texttt{Vh} are the
eigenvectors, and the eigenvalues are

\[ \lambda_k = \frac{1}{\sigma_k^2+1} \]

together with \(\lambda = 1\) repeated \(n_{\rm alm} - n_{\rm vec}\)
times for the untouched subspace. Nothing further needs computing ---
\texttt{Vh} and \texttt{Sigma} are the answer. (Checked against a dense
\texttt{eigh}: agreement to \(3\times10^{-15}\).)

These are the signal-to-noise, or Karhunen--Loève, eigenmodes familiar
from CMB analysis. Because \(\lambda_k\) \emph{is} the fraction of prior
variance surviving in mode \(k\), sorting by \(\lambda\) orders the sky
patterns from best-measured to entirely untouched. To look at one,
transform it to a map:

\begin{Shaded}
\begin{Highlighting}[]
\ImportTok{from}\NormalTok{ mistsim.alm }\ImportTok{import}\NormalTok{ alm1d\_to\_hp}
\NormalTok{mode }\OperatorTok{=}\NormalTok{ np.sqrt(Sdiag) }\OperatorTok{*}\NormalTok{ Vh[k]                       }\CommentTok{\# S\^{}\{1/2\} v\_k}
\NormalTok{m }\OperatorTok{=}\NormalTok{ hp.alm2map(np.asarray(alm1d\_to\_hp(mode)), nside)}
\end{Highlighting}
\end{Shaded}

The \(S^{1/2}\) is what puts the mode back into temperature units. Plot
\texttt{Vh{[}k{]}} directly and you get the whitened pattern, which
\(C_\ell^{-1/2}\) has tilted towards high \(\ell\).

Equivalently, \(S^{1/2}v_k\) solves the generalized problem
\(C_{\rm post}u = \lambda S u\) with the same \(\lambda_k\) --- the
alm-space statement of the same modes, orthogonal under the \(S^{-1}\)
inner product.

\hypertarget{these-are-not-the-eigenvectors-of-c_rm-post}{%
\subsubsection{\texorpdfstring{These are not the eigenvectors of
\(C_{\rm post}\)}{These are not the eigenvectors of C\_\{\textbackslash rm post\}}}\label{these-are-not-the-eigenvectors-of-c_rm-post}}

\(C_{\rm post} = S^{1/2}\tilde{C}_{\rm post}S^{1/2}\) is a
\emph{congruence}, not a similarity transform, so it does not preserve
eigenvectors. Substituting \(S^{1/2}v_k\) into
\(\mathrm{eigh}(S - BB^{\mathsf T})\) leaves a relative residual of 3.6
--- not even approximately eigenvectors.

That ordinary decomposition does exist, and is affordable (0.51 GB and
seconds at \(\ell_{\rm max} = 90\); 7.8 GB plus LAPACK workspace at
179). But it ranks modes by \emph{absolute} variance in
\(\mathrm{K}^2\), and \texttt{cl\_prior} spans \(8.1\times10^4\) across
\(\ell\) with \(C_\ell \sim \ell^{-2.1}\). Its leading modes are
therefore the monopole and dipole almost regardless of what the
instrument measured, and a mode at \(\ell = 90\) the beam never touched
still ranks near the bottom purely because its prior variance is
\(10^5\) times smaller. Use it only for an absolute error budget. For
"what does this configuration constrain", use
\(\lambda_k = 1/(\sigma_k^2+1)\).

One practical note: \texttt{save\_results} (\texttt{pipeline.py:2138})
stores \texttt{Sigma} and \texttt{nvec} but not \texttt{Vh}, so the
eigenvectors cannot be recovered from a saved \texttt{.npz} without
re-running the SVD. Full \texttt{Vh} is 93 MB per run at
\(\ell_{\rm max} = 90\); the leading 100 modes are 6.3 MB.

\begin{center}\rule{0.5\linewidth}{0.5pt}\end{center}

\hypertarget{6-sampling-a-low-rank-square-root}{%
\subsection{6. Sampling: a low-rank square
root}\label{6-sampling-a-low-rank-square-root}}

\texttt{posterior\_uncertainty} draws from the posterior without ever
forming \(C_{\rm post}\):

\begin{Shaded}
\begin{Highlighting}[]
\NormalTok{wfull }\OperatorTok{=}\NormalTok{ rng.normal(size}\OperatorTok{=}\NormalTok{(n\_alm, n\_realizations))}
\NormalTok{corr }\OperatorTok{=}\NormalTok{ Vht.T }\OperatorTok{@}\NormalTok{ (std\_reduce[:, }\VariableTok{None}\NormalTok{] }\OperatorTok{*}\NormalTok{ (Vht }\OperatorTok{@}\NormalTok{ wfull))}
\NormalTok{x\_tilde\_sim }\OperatorTok{=}\NormalTok{ wfull }\OperatorTok{{-}}\NormalTok{ corr}
\end{Highlighting}
\end{Shaded}

Written out, with \(w \sim \mathcal{N}(0, I)\):

\[ \tilde{x}_{\rm sim} = \Big[\,I - VV^{\mathsf H} + V\left(\Sigma^2+I\right)^{-1/2}V^{\mathsf H}\Big]\,w \]

\textbf{This bracket is a symmetric square root of \(C_{\rm post}\).}
The projector \(P = VV^{\mathsf H}\) and its complement \(I - P\) are
orthogonal, so the two blocks square independently:

\[ \Big[\,(I - P) + V(\Sigma^2+I)^{-1/2}V^{\mathsf H}\Big]^2 = (I-P) + V\left(\Sigma^2+I\right)^{-1}V^{\mathsf H} = C_{\rm post} \]

so \(\tilde{x}_{\rm sim}\) is an exact posterior draw. Two properties
matter:

\begin{enumerate}
\def\labelenumi{\arabic{enumi}.}
\tightlist
\item
  \textbf{Only \texttt{nvec} projections are computed}, so a draw costs
  \(O(n_{\rm alm}\cdot n_{\rm vec})\) and no dense matrix is formed.
\item
  \textbf{The unretained subspace is handled for free.} Modes outside
  \(V\) keep coefficient 1 --- full prior width --- which is exactly
  right for genuine null modes.
\end{enumerate}

\texttt{V} is real here (real packed basis), so \texttt{Vht.T} is
legitimately \(V^{\mathsf H}\).

\hypertarget{why-sample-rather-than-compute-the-variance-exactly}{%
\subsubsection{Why sample, rather than compute the variance
exactly?}\label{why-sample-rather-than-compute-the-variance-exactly}}

For the configurations this pipeline actually runs, there is no good
reason --- the exact route is both cheaper and better.

Forming \(C_{\rm post}\) densely in alm space is not prohibitive: at
\(\ell_{\rm max} = 90\), \(n_{\rm alm} = 8281\) and the matrix is 0.51
GB. What \emph{is} prohibitive is the object the plot needs --- the
pixel-space covariance \(Y C_{\rm post} Y^{\mathsf T}\), which at
\texttt{nside\ =\ 128} is \(196608^2\) entries, or 288 GB.

Neither is required. Folding the same truncated identity the other way
gives a rank-\texttt{nvec} \textbf{downdate of the prior}:

\[ \tilde{C}_{\rm post} = I - VMV^{\mathsf H}, \qquad M_k = 1 - \frac{1}{\sigma_k^2+1} = \frac{\sigma_k^2}{\sigma_k^2+1} = \sigma_k D_k \]

\[ C_{\rm post} = S - BB^{\mathsf T}, \qquad B = S^{1/2}VM^{1/2} \quad (n_{\rm alm} \times n_{\rm vec}) \]

Because the prior is isotropic,
\(\mathrm{diag}(YSY^{\mathsf T}) = \sigma^2_{\rm prior}\) --- the scalar
of section 8 --- and the whole ratio map follows in closed form:

\[ \frac{\sigma^2_{\rm post}(\hat{n})}{\sigma^2_{\rm prior}} = 1 - \frac{1}{\sigma^2_{\rm prior}}\sum_{k=1}^{n_{\rm vec}} \Big[\texttt{alm2map}(b_k)(\hat{n})\Big]^2 \]

This costs \texttt{nvec} spherical-harmonic transforms, is exact to
machine precision, and manifestly lands in \([0, 1]\) --- the data can
only ever remove variance. (Checked against a dense reference at
\(\ell_{\rm max} = 8\), \texttt{nvec\ =\ 40}, with the null space
deliberately exercised: agreement to \(2 \times 10^{-14}\).)

Measured \texttt{nvec} across the 16 executed runs is 130--1465, median
614 --- \textbf{13 of 16 fall below \texttt{n\_realizations\ =\ 1000}}.
For most runs the exact computation is therefore also the cheaper one:

\begin{longtable}[]{@{}llll@{}}
\toprule\noalign{}
& Transforms & Peak memory & Error \\
\midrule\noalign{}
\endhead
\bottomrule\noalign{}
\endlastfoot
Monte Carlo & 1000, fixed & 1.46 GB & \(\approx 4.5\,\%\) on the
variance \\
Exact downdate & \texttt{nvec}, median 614 & 93 MB & machine
precision \\
\end{longtable}

The Monte Carlo peak is dominated by \texttt{x\_sim\_map}, shape
\texttt{(1000,\ 196608)}, which is held in full before the
\texttt{np.std}.

What sampling buys that the exact route cannot:

\begin{itemize}
\tightlist
\item
  \textbf{Arbitrary nonlinear functionals.} The draws carry the full
  posterior distribution of \emph{anything} --- a sky-averaged
  temperature, band powers, the ratio between two regions --- not just
  second moments. Nothing in the current pipeline uses this.
\item
  \textbf{Cost independent of \texttt{nvec}.} When
  \(n_{\rm vec} \gg n_{\rm real}\) sampling wins; \texttt{all-nominal}
  (1465) and \texttt{all-nolake} (1296) are already in that regime.
\end{itemize}

So sampling is the more general tool, and the pipeline currently uses
none of its generality. It is best understood as the fallback for large
\texttt{nvec}, not as the default.

\begin{center}\rule{0.5\linewidth}{0.5pt}\end{center}

\hypertarget{7-truncation}{%
\subsection{7. Truncation}\label{7-truncation}}

\texttt{svds} is truncated at \(k\) modes and \texttt{select\_nvec}
(\texttt{pipeline.py:604}) then keeps \texttt{nvec} of them. Three
strategies:

\begin{itemize}
\tightlist
\item
  \texttt{"threshold"} (default) --- keep \(\sigma_k > 10^{-10}\).
  Safest: retains everything carrying measurable signal.
\item
  \texttt{"auto"} --- elbow of the singular-value curve in log space.
\item
  \texttt{"manual"} --- fixed count, with a warning if the mode at the
  cut still has Wiener factor \(D > 0.01\).
\end{itemize}

Since dropped modes are implicitly assigned the prior width, truncation
is exact for \(\sigma \to 0\) and a mild approximation for any mode cut
while \(\sigma\) is not yet \(\ll 1\) --- it slightly \emph{over}-states
the posterior variance there. That is the conservative direction, but it
is the reason the \texttt{"manual"} path warns.

\begin{center}\rule{0.5\linewidth}{0.5pt}\end{center}

\hypertarget{8-back-to-the-sky}{%
\subsection{8. Back to the sky}\label{8-back-to-the-sky}}

\begin{Shaded}
\begin{Highlighting}[]
\NormalTok{x\_sim     }\OperatorTok{=}\NormalTok{ np.sqrt(Sdiag)[:, }\VariableTok{None}\NormalTok{] }\OperatorTok{*}\NormalTok{ x\_tilde\_sim   }\CommentTok{\# un{-}whiten}
\NormalTok{x\_sim\_hp  }\OperatorTok{=}\NormalTok{ [alm1d\_to\_hp(col) }\ControlFlowTok{for}\NormalTok{ col }\KeywordTok{in}\NormalTok{ x\_sim.T]   }\CommentTok{\# packed {-}\textgreater{} healpy}
\NormalTok{x\_sim\_map }\OperatorTok{=}\NormalTok{ [hp.alm2map(xs, nside) }\ControlFlowTok{for}\NormalTok{ xs }\KeywordTok{in}\NormalTok{ x\_sim\_hp]}
\NormalTok{std\_map   }\OperatorTok{=}\NormalTok{ np.std(x\_sim\_map, axis}\OperatorTok{=}\DecValTok{0}\NormalTok{)               }\CommentTok{\# posterior sigma per pixel}
\end{Highlighting}
\end{Shaded}

Un-whitening restores the \(C_\ell\) scaling, and \texttt{alm2map}
rotates into pixel space --- where the covariance is emphatically
\emph{not} diagonal, so the per-pixel variance cannot simply be read off
\(C_{\rm post}\). The Monte Carlo recovers it from the sample scatter;
section 6 gives the exact alternative.

The denominator is the per-pixel variance of a Gaussian random field
with spectrum \(C_\ell\) (\texttt{pipeline.py:1077}):

\[ \sigma^2_{\rm prior} = \sum_\ell \frac{2\ell+1}{4\pi} C_\ell \]

A single scalar, because the prior is statistically isotropic. So the
plotted map

\[ \frac{\sigma^2_{\rm post}(\hat{n})}{\sigma^2_{\rm prior}} = \frac{\texttt{std\_map}^2}{\texttt{sigma2\_prior}} \]

is the \textbf{fraction of prior variance surviving in each pixel}: 0
means fully determined by the data, 1 means the data said nothing there.

Its structure is purely instrumental. In the 40 MHz \texttt{all-nominal}
run the ratio spans \(\approx 0.13\) in the Galactic plane to
\(\approx 0.32\) near the poles --- pixels the drift scan weights
heavily project onto high-\(\sigma_k\) modes and shrink the most.

\begin{center}\rule{0.5\linewidth}{0.5pt}\end{center}

\hypertarget{9-relation-to-the-wiener-filter-factors}{%
\subsection{9. Relation to the Wiener filter
factors}\label{9-relation-to-the-wiener-filter-factors}}

\texttt{wiener\_filter} (\texttt{pipeline.py:771}) uses a
\emph{different} factor for the posterior mean:

\begin{Shaded}
\begin{Highlighting}[]
\NormalTok{Dnum }\OperatorTok{=}\NormalTok{ Sigma[:nvec]}
\NormalTok{D }\OperatorTok{=}\NormalTok{ Dnum }\OperatorTok{/}\NormalTok{ (}\DecValTok{1} \OperatorTok{+}\NormalTok{ Dnum}\OperatorTok{**}\DecValTok{2}\NormalTok{)}
\end{Highlighting}
\end{Shaded}

which comes from the same core with one extra \(\Sigma\), contributed by
the \(\tilde{A}^{\mathsf H}\) acting on the data:

\[ C_{\rm post}\tilde{A}^{\mathsf H} = V\left(\Sigma^2+I\right)^{-1}\Sigma\, U^{\mathsf H} \]

\begin{longtable}[]{@{}llll@{}}
\toprule\noalign{}
Quantity & Factor & \(\sigma \gg 1\) & \(\sigma \ll 1\) \\
\midrule\noalign{}
\endhead
\bottomrule\noalign{}
\endlastfoot
Posterior std & \(1/\sqrt{\sigma^2+1}\) & \(\to 1/\sigma \to 0\) &
\(\to 1\) (prior width) \\
Wiener mean & \(\sigma/(1+\sigma^2)\) & \(\to 1/\sigma \to 0\) &
\(\to \sigma \to 0\) \\
\end{longtable}

Both vanish for unmeasured modes, but for opposite reasons: the
\textbf{mean} goes to zero because that is the prior mean; the
\textbf{std} goes to one because that is the prior width. Confusing the
two is the classic way to convince yourself a null mode was measured.

\begin{center}\rule{0.5\linewidth}{0.5pt}\end{center}

\hypertarget{10-monte-carlo-convergence}{%
\subsection{10. Monte Carlo
convergence}\label{10-monte-carlo-convergence}}

\texttt{std\_map} is a sample standard deviation over
\texttt{n\_realizations\ =\ 1000} draws (seeded, \texttt{seed=1420}, so
runs are reproducible). A sample std from \(N\) draws carries fractional
error

\[ \frac{\Delta\sigma}{\sigma} \approx \frac{1}{\sqrt{2(N-1)}} \approx 2.2\,\% \quad (N = 1000) \]

and the variance ratio, being a square, carries about \(4.5\,\%\).

\textbf{This is visible in the plots.} The fine-grained speckle in the
"Posterior Std Dev" and ratio panels is Monte Carlo noise, not sky
structure. It shrinks as \(1/\sqrt{N}\) --- raise
\texttt{n\_realizations} if you want those panels publication-smooth.
The large-scale plane-versus-pole gradient is real signal and is
unaffected.

\begin{center}\rule{0.5\linewidth}{0.5pt}\end{center}

\hypertarget{11-reproducing-the-figure}{%
\subsection{11. Reproducing the
figure}\label{11-reproducing-the-figure}}

\texttt{plot\_posterior\_maps} (\texttt{plotting.py:949}) draws three
Mollweide panels: posterior std, the variance ratio, and
\(\mathrm{SNR} = |{\rm map}|/\sigma\).

\begin{Shaded}
\begin{Highlighting}[]
\NormalTok{fl }\OperatorTok{=}\NormalTok{ np.ones(lmax }\OperatorTok{+} \DecValTok{1}\NormalTok{)}
\NormalTok{fl[}\DecValTok{11}\NormalTok{:] }\OperatorTok{=} \FloatTok{0.0}                                  \CommentTok{\# keep ell \textless{}= 10}
\NormalTok{best\_map }\OperatorTok{=}\NormalTok{ hp.alm2map(hp.almxfl(x\_rec, fl), nside}\OperatorTok{=}\DecValTok{128}\NormalTok{)}
\NormalTok{fig }\OperatorTok{=}\NormalTok{ msplt.plot\_posterior\_maps(std\_map, sigma2\_prior, best\_map, nside}\OperatorTok{=}\DecValTok{128}\NormalTok{)}
\end{Highlighting}
\end{Shaded}

Note the \texttt{fl} low-pass: the \emph{mean} map is truncated to
\(\ell \le 10\) for display, while \texttt{std\_map} is computed over
the full \(\ell\) range. The SNR panel is therefore a smoothed numerator
over an unsmoothed denominator --- fine for a visual check, but not a
quantity to quote.

Executed examples, all at cell 21:

\begin{itemize}
\tightlist
\item
  \texttt{notebooks/mapmaking/notebooks\_out/40mhz/all-nominal.ipynb}
\item
  \texttt{notebooks/mapmaking/notebooks\_out/25mhz/mars.ipynb}
\item
  \texttt{notebooks/mapmaking/notebooks\_out/mars-lmax\{40,60,90,179\}.ipynb}
  --- the \(\ell_{\rm max}\) series
\item
  plus the per-site combinations under \texttt{25mhz/} and
  \texttt{40mhz/}
\end{itemize}

\end{document}
